\subsection{Caso 7. ICMP: escaneo, payloads atípicos y posible túnel}
\label{caso:icmp}

\begin{situacion}
Se observa una tasa de ICMP mayor de lo habitual o paquetes con tamaños o patrones repetitivos. ICMP se especifica en el RFC~792 \parencite{rfc792}.
\end{situacion}

\paragraph{Procedimiento.}
\begin{itemize}[leftmargin=*, topsep=2pt, itemsep=2pt]
	\item En Wireshark, filtrar \texttt{icmp} y agrupar por IP de origen/destino, tipo, código, longitud y tiempo.
	\item Establecer una línea base de ping de diagnóstico legítimo y compararla con la captura analizada.
	\item Usar PCAP sintético o didáctico para estudiar payloads; no transmitir información sensible por ICMP.
\end{itemize}

\paragraph{Indicadores que deben buscarse.}
Ráfagas de echo request/reply, longitudes inusuales, periodicidad repetitiva o volumen no explicado por tareas de diagnóstico.

\paragraph{Mitigación y defensa.}
Permitir solo los tipos ICMP necesarios para la operación y el diagnóstico; aplicar límites proporcionados; monitorizar sin bloquear indiscriminadamente mensajes requeridos por la red.

\paragraph{Evidencias que se entregan.}
Dashboard por IP, filtro aplicado, tabla de tamaños y frecuencias, y propuesta de regla.

\paragraph{Preguntas de análisis.}
\begin{enumerate}[leftmargin=*, topsep=2pt, itemsep=2pt]
	\item ¿Qué funciones legítimas cumple ICMP?
	\item ¿Qué impacto tendría bloquearlo completamente?
\end{enumerate}

\begin{table}[H]
	\caption{Distribución de tráfico ICMP por tipo y longitud}
	\label{tab:icmp}
	\centering
	\footnotesize
	\renewcommand{\arraystretch}{1.25}
	\begin{tabularx}{\textwidth}{@{}c c c >{\raggedright\arraybackslash}p{0.17\textwidth} >{\raggedright\arraybackslash}p{0.17\textwidth} Y@{}}
		\toprule
		\textbf{Tipo} & \textbf{Código} & \textbf{Longitud} & \textbf{Origen} & \textbf{Destino} & \textbf{Frecuencia / observación} \\
		\midrule
		& & & & & \\
		& & & & & \\
		& & & & & \\
		& & & & & \\
		& & & & & \\
		\bottomrule
	\end{tabularx}
\end{table}

\registrohallazgo{Caso 7}
